%% notes-tex/026-kinetics-thermodynamics/sections/5-module.tex
%% [[torchcell.models.equivariant_cell_graph_transformer.mermaid.metabolic-module]]
%% https://github.com/Mjvolk3/torchcell/tree/main/notes-tex/026-kinetics-thermodynamics/sections/5-module.tex

\section{Where these tables enter the model\secstatus{todo}}

Figure~\ref{fig:cgt-module} places the flux module inside the cell graph transformer. The
module is neither of the two instrument classes the architecture previously had: a Type I
instrument maps a representation to a representation and a Type II instrument maps a
representation to a measured output, while the flux module maps a representation to a
\emph{flux vector}, a latent physical state from which phenotypes are then read. It sits
between them.

Three of its inputs are parameter tables rather than learned weights, and those three are
what this document builds. The $k_{cat}$ table sets the upper bound of the box that the
flux is parameterized inside, so it fixes what fluxes are reachable at all. The formation
energies set the reaction potential whose sign the second-law penalty acts on, which is
what makes that penalty a physical statement rather than a regularizer. The $K_M$ table is
built and not yet enabled, for the reason Section~\ref{sec:limits} gives.

The heterologous cassette enters as a typed perturbation on the genome-scale model, which
is what puts a betaxanthin pathway and a gene deletion inside the same object rather than
in two parallel code paths.

\tcfigfit[1.0]{figures/cgt_metabolic_module.pdf}{The cell graph transformer with the
  enzyme-constrained thermodynamic flux module. Purple marks inputs and the hash-pinned
  parameter tables built here; orange the embeddings and the Type I perturbation operators;
  red the transformer layers and the outputs; yellow the graph regularization, the
  invariant readouts and the soft constraints; blue the metabolic module, whose color slot
  was unused in the previous version of this diagram so the new content separates at a
  glance. Every hard constraint is either exact by construction, as the box and the
  capacity bound are, or relaxed to a smooth penalty; nothing is enforced by a binary
  variable, which is what makes the layer differentiable and what removes the mixed-integer
  cost a per-genotype thermodynamic program pays. The dashed $K_M$ edge marks a table that
  exists but is not yet consumed. Scaled to the text block: the source renders at 211.7 mm
  wide against a 182 mm block, so this is not shown at true print size. Source:
  \texttt{notes/torchcell.models.equivariant\_cell\_graph\_transformer.mermaid.metabolic-module.md},
  rendered by \texttt{notes/assets/publish/scripts/mermaid\_pdf.sh}.}{fig:cgt-module}
